\par\noindent\textbf{(1)} 晶格中的半经典动力学满足ℏq̇=F与 ẋ=(1/ℏ)∂E/∂q。故 q(t)=q₀+Ft/ℏ，经过一个倒格矢2π/a需时间T\_B=2πℏ/(|F|a)，Bloch角频率 ω\_B=|F|a/ℏ。q在约化布里渊区循环而非真实动量无限增加。
\par\noindent\textbf{(2)} 能带导数为 ∂E/∂q=2Ja sin(qa)，因此 ẋ=(2Ja/ℏ)sin(q₀a+Fat/ℏ)。积分得到 x(t)−x(0)=−(2J/F)[cos(q₀a+Fat/ℏ)−cos(q₀a)]，完全周期且平均净位移为零。
\par\noindent\textbf{(3)} 相对振荡中心的位移幅值为2J/|F|，峰谷跨度为4J/|F|。J→0时能带扁平、群速度与运动幅度均归零；F→0时周期发散，应先对cos差作小F展开，得到匀速传播x≈(2Ja/ℏ)sin(q₀a)t，而不能将幅值公式误判为无限物理振荡。
\par\noindent\textbf{(4)} 若区边界与下一能带的最小间隙为Δ\_gap，力使准动量以速率F/ℏ穿越避交叉，两能级Landau–Zener隧穿几率量级 P\_LZ≈exp[−πΔ\_gap²/(2ℏ|d(ε₁−ε₂)/dt|)]（此处Δ\_gap定义为最小完整能级间距）。当扫掠足够慢、Δ\_gap²大于ℏ扫掠能量速率时带间泄漏指数抑制；大F时单带Bloch振荡不再封闭。
